\documentclass[12pt,a4paper]{article}
\usepackage{amsmath,amssymb,amsthm}
\newtheorem{proposition}{Proposition}[section]
\newtheorem{corollary}[proposition]{Corollary}
\theoremstyle{remark}
\newtheorem*{remark}{Remark}
\begin{document}
\section{Finite groups}
Let $G$ be a finite group of order $n$. Recall the following facts.
\begin{enumerate}
\item Every subgroup $H\le G$ has order dividing $n$.
\item For each $g\in G$ the order of $g$ divides $n$.
\item If $n$ is prime, then $G$ is cyclic:
  \begin{enumerate}
  \item pick any $g\ne e$;
  \item the subgroup $\langle g\rangle$ has order greater than $1$ dividing $n$;
  \item hence $\langle g\rangle=G$.
  \end{enumerate}
\end{enumerate}

\begin{proposition}\label{prop:lagrange}
If $|G|=n$ and $g\in G$, then $g^n=e$.
\end{proposition}
\begin{proof}
The order $d$ of $g$ divides $n$ by item (2) above, say $n=dk$. Thus $g^n=(g^d)^k=e$.
\end{proof}

\begin{corollary}
In the multiplicative group $(\mathbb{Z}/p\mathbb{Z})^\times$ we have $a^{p-1}\equiv1\pmod p$ for every $a\not\equiv0$.
\end{corollary}

\begin{remark}
Proposition~\ref{prop:lagrange} is Lagrange's theorem in disguise; the corollary is Fermat's little theorem.
\end{remark}

\section{Examples}
\begin{description}
\item[Cyclic] $\mathbb{Z}/n\mathbb{Z}$ under addition.
\item[Symmetric] $S_n$, of order $n!$.
\item[Dihedral] $D_n$, of order $2n$, generated by a rotation $r$ and a reflection $s$ with $srs=r^{-1}$.
\end{description}
\end{document}
